\section{总结与延伸}

杨植麟的演讲构建了一条从\textbf{第一性原理到工程实践}的清晰技术路线：

\begin{enumerate}
\item \textbf{第一性原理}：Scaling Law——能源转化为智能
\item \textbf{两个优化维度}：Token Efficiency $\times$ Long Context
\item \textbf{优化器革命}：Muon二阶优化器（2倍Token Efficiency）+ QK Clip（训练稳定性）
\item \textbf{架构创新}：Kimi Linear / Delta Attention（线性注意力首次全面超越全注意力）
\item \textbf{产品落地}：Kimi K2（中国首个Agentic模型，HLE超越OpenAI）
\item \textbf{未来方向}：K3（基于Kimi Linear） $\rightarrow$ K4 $\rightarrow$ K5 $\rightarrow$ K100
\end{enumerate}

\begin{knowledgebox}{拓展阅读}
\begin{itemize}
\item Kaplan et al., ``Scaling Laws for Neural Language Models'' (2020)：最早系统性阐述Scaling Law
\item Muon优化器相关工作：二阶优化在大模型预训练中的应用
\item Kimi Linear技术报告：Delta Attention的形式化推导与工程实现
\item HLE (Humanities Last Exam)：Benchmark说明
\end{itemize}
\end{knowledgebox}

%% --- Fin del contenido --- %%

\end{document}
